\begin{corollary}[6.13.5]
\label{IV.2.s6.13.5-fr}
Si $X$ es tal que $\operatorname{Reg}(X)$ es abierto en $X$, entonces
$\operatorname{Nor}(X)$ es abierto en $X$.
\end{corollary}

\begin{proof}
En efecto, todo punto maximal $x$ en el que $X$ es reducido (y en el
que, por tanto, $\mathcal O_{X,x}$ es un cuerpo) pertenece a
$\operatorname{Reg}(X)$; por hipótesis, es entonces un punto interior
de $\operatorname{Reg}(X)$ y, \emph{con mayor razón}, un punto interior de
$\operatorname{Nor}(X)$.
\end{proof}
